% Vietnamese translation for OI Public Library
% Source-PDF: USACO/【由 BirDOR 整理】USACO 2001-2007 月赛测试数据+题目+题解/USACO月赛十年题典v10.pdf
\documentclass[11pt]{article}
\usepackage{oipl-vn}
\usepackage{array}
\usepackage{booktabs}
\usepackage{longtable}

\title{USACO: Tuyển tập đề thi tháng trong mười năm, bản v10}
\author{}
\date{}

\newcolumntype{L}[1]{>{\raggedright\arraybackslash}p{#1}}
\newcommand{\codes}[1]{{\footnotesize\ttfamily #1}}
\newcommand{\monthrow}[3]{#1 & #2 & \codes{#3}\\}
\newcommand{\problem}[4]{%
  \subsection[#1 (#2)]{#1 \hfill {\normalsize\texttt{#2}}}%
  \noindent\textit{#3. #4}\par\smallskip
}
\newcommand{\iosec}[1]{\paragraph{#1}}

\begin{document}
\maketitle

\origtitle{USACO Monthly Contest Problem Anthology in Ten Years, v10}
\sourcepath{USACO/.../USACO-monthly-contest-ten-year-v10.pdf}

\section*{Ghi chú về bản dịch}

PDF nguồn dài 266 trang, tạo bằng LaTeX, có tiêu đề ``USACO Monthly
Contest Problem Anthology in Ten Years'' trong metadata. Trang bìa ghi
USACO, tuyển tập đề thi tháng trong mười năm, tác giả Zhuang Le, địa chỉ
liên hệ \texttt{birdor@126.com}, ngày 27/4/2010.

Văn bản trích bằng \texttt{pdftotext} bị lỗi mã hóa với phần tiếng Trung:
các tiêu đề, đoạn văn và phát biểu bài phần lớn trở thành mojibake. Tôi đã
render các trang mở đầu bằng \texttt{pdftoppm} để đọc trực quan; nhờ vậy có
thể khôi phục trang bìa, lời tựa và khung mục lục. Thân tài liệu không được
dịch toàn văn trong ghi chú này; phần có thể khôi phục ổn định là năm thi,
đợt thi và mã bài trong ngoặc vuông. Các mã bài dưới đây giữ nguyên để người
đọc đối chiếu với kho USACO hoặc các bản dịch chi tiết khác trong thư mục
USACO của bản tiếng Việt.

\section{Lời tựa}

\subsection{Lý do biên soạn}

Trong thế giới OI, từ NOIP đến ACM, số lượng bài tập rất lớn. Nhiều thí sinh
khi gặp một biển đề như vậy không biết nên bắt đầu từ đâu. USACO có uy tín
tốt, cách ra đề khá gần với NOI và các kỳ chọn đội cấp tỉnh, nên rất phù hợp
để luyện tập thường xuyên.

Tác giả chia bài USACO thành hai nhóm: bài huấn luyện và bài thi tháng. Bài
huấn luyện đã được nhiều thế hệ học sinh dịch, sắp xếp và biết đến rộng rãi;
ngược lại, nhiều bài thi tháng vẫn chưa được dịch và chưa được hệ thống hóa.
Điều đó khiến người học gặp rào cản khi muốn dùng các đề thi tháng để nâng
trình. Vì vậy tác giả thực hiện công việc dịch và sắp xếp tuyển tập này.

\subsection{Vấn đề bản quyền}

Các đề gốc trong tuyển tập đến từ trang thi tháng USACO
\texttt{http://contest.usaco.org}; tên tác giả gốc được đặt ở mục thông tin
nguồn của từng bài. Bản dịch tiếng Trung có nhiều nguồn: một phần do trang
USACO mời tuyển thủ Trung Quốc dịch, một phần từ diễn đàn, một phần từ các
tiền bối của tác giả, và phần còn lại do tác giả tự dịch. Tài liệu không dùng
cho mục đích thương mại và không phải ấn phẩm xuất bản chính thức. Bản quyền
đề gốc thuộc USACO; bản dịch thuộc người dịch tương ứng.

\subsection{Cách sử dụng tuyển tập}

Tác giả nhắc rằng những bản cũ có thể chứa lỗi ảnh hưởng đến lời giải, nên
nên dùng bản mới nhất nếu có. Với học sinh luyện thi NOIP, tài liệu khuyến
nghị tập trung vào các bài mức đồng và bạc, cùng một số bài xanh sớm. Do bài
mức vàng không hoàn toàn khớp với hướng ra đề của NOIP, có thể đọc tham khảo
nhưng không cần đi quá sâu. Với học sinh hướng tới tuyển chọn đội tỉnh hoặc
NOI, tài liệu khuyến nghị tập trung vào mức xanh và vàng, đồng thời làm thêm
một lượng vừa phải bài đồng và bạc.

Một số nhóm bài không được thu thập đầy đủ: bài đồng trước năm 2003, bài xanh
năm 2004 và các bài đồng về sau được xem là quá cơ bản; một số bài tháng mười
về sau cũng không được đưa vào vì lý do tương tự. Phần lớn bài nộp theo mẫu
và bài tương tác không được thu thập. Với các chủ đề như tìm kiếm và luồng
mạng, tác giả khuyên người học bổ sung bằng nguồn khác. Sau khi làm xong,
người học nên kiểm thử chương trình, có thể nộp lên kho bài ACM của Đại học
Bắc Kinh hoặc tải dữ liệu từ trang USACO chính thức để tự kiểm thử.

\subsection{Lời mời góp ý}

Do năng lực có hạn, tuyển tập chắc chắn còn sai sót và cần cập nhật. Tác giả
mời người đọc gửi thư chỉ ra các lỗi như bản dịch không khớp đề gốc, dùng từ
không đúng, ví dụ sai, lỗi chính tả, lỗi dấu câu hoặc lỗi dàn trang. Khi góp
ý, người đọc nên ghi rõ số phiên bản đã tham khảo.

\section{Mục lục khôi phục}

Bảng sau là cấu trúc mục lục có thể đọc lại ổn định từ PDF: mỗi dòng ghi năm,
đợt thi và các mã bài xuất hiện trong đợt đó. Tên tiếng Trung của từng bài
không được đưa vào vì phần trích văn bản bị lỗi mã hóa và yêu cầu bản dịch
không để lẫn chữ Hán trong nội dung hiển thị.

\begin{longtable}{L{0.12\linewidth}L{0.18\linewidth}L{0.62\linewidth}}
\toprule
\textbf{Năm} & \textbf{Đợt thi} & \textbf{Mã bài khôi phục}\\
\midrule
\endfirsthead
\toprule
\textbf{Năm} & \textbf{Đợt thi} & \textbf{Mã bài khôi phục}\\
\midrule
\endhead
\bottomrule
\endlastfoot
2000 & Spring & \codes{ski farm route mat}\\
2000 & Open & \codes{digit maze evac knight chore}\\
2000 & Fall & \codes{outfrnd infrnd enemy amicbl}\\
\midrule
2001 & Spring & \codes{spots cowfix route barn}\\
2001 & Open & \codes{relay quake sort sign well}\\
2001 & Fall & \codes{cowdom plumb dinner prefix}\\
2001 & Winter & \codes{split count treas pool bed1}\\
\midrule
2002 & February & \codes{fiber power cycling roads pasture}\\
2002 & Spring & \codes{hypno happy ucc}\\
2002 & Open & \codes{majesty secret tix life}\\
2002 & Fall & \codes{palpath apples therock}\\
2002 & Winter & \codes{pasture party stroll grazeset}\\
2002 & December & \codes{nextree journal captives btp}\\
\midrule
2003 & February & \codes{cowmath tour impster traffic}\\
2003 & March & \codes{twofour cowfnc cornfld}\\
2003 & Open & \codes{mtwalk leap2 milking}\\
2003 & November & \codes{smrtfun mgrid popular msworld}\\
2003 & December & \codes{cover elite sprbrk cowq}\\
\midrule
2004 & January & \codes{order lock arithprg flood rects sixdeg}\\
2004 & February & \codes{navigate marathon dquery dstats breeding thegame parens}\\
2004 & March & \codes{tryout pizza finance liecow serial paranoid}\\
2004 & Open & \codes{cubes lineup moofest turnin cavecow1 cavecow2 cavecow3 cavecow4}\\
2004 & November & \codes{bcatch lkcount cowhome middle bullmath bankint}\\
2004 & December & \codes{divide obstacle skiarea cleaning cowtract treecut}\\
\midrule
2005 & January & \codes{cover juice naptime sumset watchcow volume}\\
2005 & February & \codes{jpoi secret aggr acquire cowrig fcount}\\
2005 & March & \codes{ombro elevator yogfac alibi outofhay}\\
2005 & Open & \codes{lazy exp around peaks waves navcit disease mud}\\
2005 & October & \codes{cowski flight nearfr allow}\\
2005 & November & \codes{asteroid ontherun twalk skyline acrobat ants}\\
2005 & December & \codes{cpattern expand layout ni clean scales}\\
\midrule
2006 & January & \codes{rpaths roping corral prom ddayz grove}\\
2006 & February & \codes{trt bdsum cell stead reserve}\\
2006 & March & \codes{skilift tselect mooo slides xlite}\\
2006 & Open & \codes{cfair mqueue twohead events wall}\\
2006 & October & \codes{acng skate pie1 lineup moat}\\
2006 & November & \codes{plank cowfood block badhair bigsq rndnum}\\
2006 & December & \codes{wormhole fewcoins patterns picnic coaster jump}\\
\midrule
2007 & January & \codes{lineupg psolve schul flowers tallest}\\
2007 & February & \codes{newbarn csort lilypad lexicon silvlily sparty}\\
2007 & March & \codes{lineup ranking cowturn traffic balance expense}\\
2007 & Open & \codes{cheappal dining horizon catchcow fliptile}\\
2007 & October & \codes{money paint2 secpas obstacle telewire relays tanning hurdles milkprod bcl}\\
2007 & Special Chinese & \codes{sumsums baabo treasure}\\
2007 & December & \codes{sightsee gourmet bclgold charm roads mud}\\
\midrule
2008 & January & \codes{bales tower contest cowrun phoneline}\\
2008 & February & \codes{grading hotel lines meteor egroup}\\
2008 & March & \codes{acquire cowjog ppairing baler ctravel river}\\
2008 & Open & \codes{crisis nabor ratf wordpow cowcar danger}\\
2008 & October & \codes{pwrfail bones water quad pwalk rotation}\\
2008 & November & \codes{mixup2 cheer lites toy buyhay guardian mtime}\\
2008 & December & \codes{treat sec winchk fence hay4sale patheads jigsaw}\\
\midrule
2009 & Holiday & \codes{holpaint cattleb delay}\\
2009 & January & \codes{damage baric travel lphone bestspot flow}\\
2009 & February & \codes{shuttle stock revamp lepr surround bullcow}\\
2009 & March & \codes{baying cleanup damage2 sandcas fristeam lookup}\\
2009 & Open & \codes{ski job tower cowemb hideseek cline cdgame graze2}\\
2009 & October & \codes{milkweed allow diet heatwv}\\
2009 & November & \codes{lights cookie rescue xoinc farmoff jobhunt}\\
2009 & December & \codes{dizzy toll vidgame sgraze bobsled mnotes}\\
\end{longtable}

\section{Các bài đã khôi phục từ OCR}

Phần dưới đây bắt đầu thay thế ghi chú chỉ-mục bằng các phát biểu bài đã đọc
lại từ ảnh render của PDF và OCR bằng \texttt{tesseract -l chi\_sim+eng}. Mã
bài, tên tệp, dữ liệu mẫu và tên tác giả nguồn được giữ nguyên để đối chiếu.
Những khối mã chương trình, nếu có ở các trang sau, vẫn thuộc ngoại lệ không
dịch code template của dự án.

\section{USACO 2000 Spring}

\problem{Trượt tuyết}{ski}{Nhóm xanh}{Brian Dean}

\iosec{Tóm tắt bài toán.}
Để thưởng cho đàn bò, Farmer John muốn đưa chúng tới Colorado thăm huấn luyện
viên Rob của USACO và đi trượt tuyết. Một khu trượt tuyết được mô tả bằng các
điểm trượt, độ cao của từng điểm, và các đường nối giữa các điểm. Bò chỉ có
thể trượt từ điểm cao hơn xuống điểm thấp hơn qua một đường nối có sẵn.

Điểm cao nhất là đỉnh núi, điểm thấp nhất là chân núi. Rob muốn đưa càng nhiều
bò càng tốt, nhưng hai con bò bất kỳ không được đi đúng cùng một tuyến; chỉ
cần có một cạnh trên đường đi khác nhau thì hai tuyến được xem là khác nhau.
Hãy tính số bò nhiều nhất có thể đưa đi trượt.

\iosec{Dữ liệu vào.}
Dòng đầu gồm \(N,M\), với \(2<N<200\) và \(1<M<5000\), lần lượt là số điểm
trượt và số đường nối. Tiếp theo là \(N\) số cho biết độ cao của các điểm,
mỗi độ cao không quá \(1000000\). Sau đó là \(M\) cặp số, mỗi cặp mô tả một
đường nối giữa hai điểm. Không có cạnh tự nối, nhưng giữa hai điểm khác nhau
có thể có nhiều đường nối.

\iosec{Dữ liệu ra.}
In một số nguyên: số bò nhiều nhất Rob có thể đưa đi trượt.

\iosec{Ví dụ.}
\begin{verbatim}
4 5
500 400 300 200
1 2 2 3 3 4 1 4 2 4
\end{verbatim}

\begin{verbatim}
3
\end{verbatim}

\problem{Nông trại}{farm}{Nhóm xanh}{Brian Dean; dịch giả nguồn: Liu Yu}

\iosec{Tóm tắt bài toán.}
Farmer John muốn vẽ bản đồ nông trại của mình. Bản đồ là một hình vuông
\(1\times 1\) km được chia thành các vùng chữ nhật nhỏ, mỗi vùng có một màu.
Mục tiêu ban đầu của John là thiết kế sao cho hai vùng kề nhau bất kỳ có màu
khác nhau.

Mô tả của nông trại được cho dưới dạng đệ quy. Một vùng có thể được tách bởi
một đường thẳng đứng, bởi một đường thẳng ngang, hoặc không tách nữa và khi đó
được gán màu. Hãy xác định có bao nhiêu cặp vùng kề nhau có cùng màu.

\iosec{Dữ liệu vào.}
Đầu vào là một mô tả gồm các số màu và hai từ khóa \texttt{vsplit},
\texttt{hsplit}. \texttt{vsplit} tách vùng hiện tại thành phần trái và phần
phải; \texttt{hsplit} tách vùng hiện tại thành phần trên và phần dưới. Nếu
vùng không bị tách nữa, dòng hiện tại là màu của vùng. Tổng số vùng không quá
\(100\).

\iosec{Dữ liệu ra.}
In một số nguyên: số cặp vùng kề nhau có cùng màu.

\iosec{Ví dụ.}
\begin{verbatim}
vsplit
5
hsplit
vsplit
37
6
5
\end{verbatim}

\begin{verbatim}
2
\end{verbatim}

\problem{Tuyến đường của bò}{route}{Nhóm xanh}{Brian Dean; dịch giả nguồn: Liu Yu}

\iosec{Tóm tắt bài toán.}
Một xe thư đi ngang qua Farmer John đâm vào tảng đá, làm rơi ra nhiều bưu
thiếp. Mỗi bưu thiếp mô tả một tuyến đường từ một thành phố tới một thành phố
khác, bằng cách ghi thành phố xuất phát rồi lần lượt ghi hướng đi, khoảng cách
và thành phố kế tiếp. Các hướng có thể là đông, nam, tây, bắc.

Hãy kiểm tra các mô tả tuyến đường có nhất quán với nhau hay không, tức là có
thể đặt các thành phố trên cùng một mặt phẳng sao cho mọi mô tả đều đúng hay
không. Kết quả cần tìm là tiền tố dài nhất của danh sách mô tả vẫn còn hợp lệ.

\iosec{Dữ liệu vào.}
Dòng đầu là số tuyến \(R\), với \(2\le R\le 200\). Mỗi mô tả bắt đầu bằng thành
phố xuất phát, sau đó là một dãy bộ ba: hướng \texttt{N}, \texttt{S},
\texttt{E}, \texttt{W}; khoảng cách trong \([1,200]\); và thành phố kế tiếp.
Thành phố được đánh số trong \([1,200]\). Số \(0\) kết thúc mô tả hiện tại.
Một mô tả có thể trải trên nhiều dòng, nhưng mô tả mới luôn bắt đầu ở một dòng
mới; trong một tuyến, không thành phố nào xuất hiện hai lần.

\iosec{Dữ liệu ra.}
In một số nguyên \(n\): số mô tả đầu tiên có thể đồng thời đúng, với \(n\) lớn
nhất có thể.

\iosec{Ví dụ.}
\begin{verbatim}
2
1 3 N 2 2 W 3 1 S 4 0
4 3 W 1 0
\end{verbatim}

\begin{verbatim}
1
\end{verbatim}

\problem{Ma trận lớn nhất}{mat}{Nhóm xanh}{Adapted from 1995 UMD Competition; dịch giả nguồn: Liu Yu}

\iosec{Tóm tắt bài toán.}
Một ma trận tăng là ma trận mà mỗi hàng và mỗi cột đều tăng. Cho một ma trận
có các phần tử nằm trong \([0,32000]\), hãy tìm mọi ma trận con tăng có diện
tích lớn nhất. Nếu có nhiều ma trận con cùng đạt diện tích lớn nhất, in chúng
theo thứ tự tự nhiên của dãy đầu ra: so sánh số thứ nhất, rồi số thứ hai, v.v.

\iosec{Dữ liệu vào.}
Dòng đầu có hai số \(r,c\), với \(1<r,c<150\), là số hàng và số cột. Sau đó có
\(r\cdot c\) số, theo thứ tự từ trái sang phải, từ trên xuống dưới. Vì số lượng
giá trị trên mỗi dòng không cố định, chương trình gốc khuyến nghị đọc bằng
kiểu đọc bỏ qua xuống dòng.

\iosec{Dữ liệu ra.}
In tất cả ma trận con tăng lớn nhất, mỗi ma trận con trên một dòng. Hai số đầu
là tọa độ hàng và cột của góc trái trên; hai số sau là chiều rộng và chiều cao
của ma trận con.

\iosec{Ví dụ.}
\begin{verbatim}
3 4
8 2 3 9
3 5 7 8
7 2 1 9
\end{verbatim}

\begin{verbatim}
1 2 2 2
2 1 1 4
\end{verbatim}

\section{USACO 2000 Open}

\problem{Tiếp sức chữ số}{digit}{Nhóm xanh}{Russ Cox; dịch giả nguồn: Liu Yu}

\iosec{Tóm tắt bài toán.}
Đàn bò chơi một trò chơi với số nguyên. Bắt đầu từ một số như \(6593\), lấy
tất cả chữ số của nó nhân với nhau để được số mới: \(6\cdot5\cdot9\cdot3=810\).
Lặp lại quá trình này cho tới khi thu được một số có một chữ số. Hãy in toàn
bộ dãy số sinh ra trong quá trình chơi.

\iosec{Dữ liệu vào.}
Một số nguyên \(N\), với \(10<N<2\cdot 10^9\).

\iosec{Dữ liệu ra.}
Trên một dòng, in theo thứ tự mọi số xuất hiện trong trò chơi, kết thúc ở số
có một chữ số.

\iosec{Ví dụ.}
\begin{verbatim}
98886
\end{verbatim}

\begin{verbatim}
98886 27648 2688 768 336 54 20 0
\end{verbatim}

\problem{Sơ tán an toàn}{evac}{Nhóm xanh}{Eljakim Schrijvers; dịch giả nguồn: Liu Yu}

\iosec{Tóm tắt bài toán.}
Farmer John muốn bảo đảm rằng trong tình huống khẩn cấp, mọi con bò đều có một
phương án thoát thân an toàn. Khi hoảng loạn, bò chỉ chạy thẳng theo hướng mà
John đã dặn trước: hoặc lên phía bắc, tức giảm chỉ số hàng đi \(1\), hoặc sang
phía đông, tức tăng chỉ số cột đi \(1\).

Để tránh va chạm, đường thoát của bất kỳ con bò nào không được đi qua vị trí
ban đầu của một con bò khác. Cho các vị trí bò trên một lưới \(r\) hàng,
\(c\) cột, hãy xác định bản đồ đã an toàn chưa; nếu chưa, liệu bỏ đúng một con
bò có làm bản đồ an toàn hay không.

\iosec{Dữ liệu vào.}
Dòng đầu gồm \(r,c\), số hàng và số cột. Dòng thứ hai là \(n\), số bò. \(n\)
dòng tiếp theo, mỗi dòng gồm hàng và cột của một con bò.

\iosec{Dữ liệu ra.}
Nếu bản đồ đã an toàn, in \texttt{0}. Nếu bỏ bất kỳ một con bò nào mà bản đồ
vẫn không an toàn, in \texttt{-1}. Ngược lại, in \texttt{1}, rồi ở dòng tiếp
theo in chỉ số con bò cần bỏ; nếu có nhiều lựa chọn, lấy chỉ số nhỏ nhất theo
thứ tự nhập.

\iosec{Ví dụ.}
\begin{verbatim}
5 5
5
1 1
2 4
3 1
2 2
2 1
\end{verbatim}

\begin{verbatim}
1
5
\end{verbatim}

\problem{Quân mã}{knight}{Nhóm xanh}{Nguồn không ghi rõ trong trang OCR}

\iosec{Tóm tắt bài toán.}
Đàn bò chuyển từ các biến thể của bài toán \(n\)-hậu sang bài toán quân mã.
Một nước đi được ký hiệu \([a,b]\): nếu \(a>0\) thì đi lên \(a\) bước, nếu
\(a<0\) thì đi xuống; tương tự, \(b>0\) là sang phải và \(b<0\) là sang trái.

Với quy tắc \((A,B)\), quân mã có tám nước đi:
\[
[+A,+B],\ [+A,-B],\ [-A,+B],\ [-A,-B],
\]
\[
[+B,+A],\ [+B,-A],\ [-B,+A],\ [-B,-A].
\]
Quân mã thường có quy tắc \((1,2)\). Cho bàn cờ \(N\times N\), hãy tìm số
quân mã ít nhất cần đặt để khống chế toàn bộ bàn cờ. Một quân mã khống chế các
ô nó đi tới được trong một nước, nhưng không khống chế chính ô nó đứng.

\iosec{Dữ liệu vào.}
Dòng đầu là \(N\), với \(4<N<8\). Dòng thứ hai gồm \(A,B\), mô tả quy tắc đi.

\iosec{Dữ liệu ra.}
In số quân mã ít nhất cần đặt.

\iosec{Ví dụ.}
\begin{verbatim}
6
1 2
\end{verbatim}

\begin{verbatim}
8
\end{verbatim}

\problem{Việc vặt}{chore}{Nhóm xanh}{Don Piele; dịch giả nguồn: XY}

\iosec{Tóm tắt bài toán.}
Trước khi vắt sữa, nông trại của Farmer John có nhiều việc phải hoàn thành:
tập hợp bò, lùa bò vào chuồng, rửa bầu vú, v.v. Một số việc chỉ được bắt đầu
sau khi các việc chuẩn bị của nó đã hoàn thành.

John có \(n\) việc, mỗi việc có thời gian thực hiện riêng. Danh sách được sắp
theo thứ tự phụ thuộc: với việc \(k>1\), các việc chuẩn bị của nó chỉ có thể
nằm trong các việc \(1,\ldots,k-1\). Các việc không phụ thuộc nhau có thể làm
đồng thời, và nông trại có đủ công nhân. Hãy tính thời gian ngắn nhất để hoàn
thành tất cả việc.

\iosec{Dữ liệu vào.}
Dòng đầu là \(n\), với \(3<n<10000\). Sau đó có \(n\) dòng, mỗi dòng mô tả một
việc: số thứ tự việc, thời gian \(\mathit{len}\) với \(1<\mathit{len}<100\),
rồi danh sách các việc chuẩn bị, kết thúc bằng \(0\). Tổng số quan hệ chuẩn bị
không quá \(100\).

\iosec{Dữ liệu ra.}
In thời gian ngắn nhất để hoàn thành mọi việc.

\iosec{Ví dụ.}
\begin{verbatim}
7
1 5 0
2 2 1 0
3 3 2 0
4 6 1 0
5 1 2 4 0
6 8 2 4 0
7 4 3 5 6 0
\end{verbatim}

\begin{verbatim}
23
\end{verbatim}

\section{USACO 2000 Fall}

\problem{Bạn bò ngoài trời}{outfrnd}{Nhóm xanh}{Brian Dean, 2000; dịch giả nguồn: BirDOR}

\iosec{Tóm tắt bài toán.}
John có \(N\) con bò đánh số từ \(1\) đến \(N\), và \(K\) bãi cỏ xếp thành một
hàng. Mỗi sáng, bò đi theo thứ tự số hiệu; tại mỗi bãi cỏ, John giữ lại một số
con đầu hàng còn lại, rồi các con sau tiếp tục đi tới bãi kế tiếp. Mỗi bò sau
khi vào bãi cỏ sẽ ở đó cả ngày.

Một cặp bạn bò \((i,j)\) có một mong muốn được ở cùng bãi. Hãy chia dãy bò
thành \(K\) đoạn liên tiếp để tối đa hóa số mong muốn được thỏa mãn. Không bãi
cỏ nào được rỗng, và không bãi nào được chỉ có một con bò.

\iosec{Dữ liệu vào.}
Dòng đầu gồm \(N,K,F\), với \(4\le 2K<N<150\) và \(1<F<200\), trong đó \(F\)
là số cặp bạn. \(F\) dòng tiếp theo, mỗi dòng là một cặp bò bạn bè.

\iosec{Dữ liệu ra.}
In số mong muốn lớn nhất có thể thỏa mãn.

\iosec{Ví dụ.}
\begin{verbatim}
5 2 4
1 2
2 3
4 5
1 5
\end{verbatim}

\begin{verbatim}
3
\end{verbatim}

\problem{Bạn bò trong nhà}{infrnd}{Nhóm xanh}{Hal Burch, 2000; dịch giả nguồn: BirDOR}

\iosec{Tóm tắt bài toán.}
Buổi tối, John phải xếp giường cho \(N\) con bò. Có đúng \(N\) giường, đánh số
từ \(1\) đến \(N\). Với mỗi cặp bò bạn bè, khoảng cách của cặp đó là trị tuyệt
đối của hiệu số giường mà hai con nằm. Hãy xếp bò vào giường sao cho tổng
khoảng cách của mọi cặp bạn bè là nhỏ nhất.

\iosec{Dữ liệu vào.}
Dòng đầu gồm \(N,F\), với \(3<N<30\) và \(1<F<4N\). \(F\) dòng tiếp theo, mỗi
dòng là một cặp bò bạn bè.

\iosec{Dữ liệu ra.}
In tổng khoảng cách nhỏ nhất.

\iosec{Ví dụ.}
\begin{verbatim}
6 5
1 2
1 3
2 5
3 6
5 6
\end{verbatim}

\begin{verbatim}
8
\end{verbatim}

Một cách xếp đạt được là:
\begin{verbatim}
1, 2, 3, 5, 6, 4.
\end{verbatim}

\problem{Kẻ thù của nhau}{enemy}{Nhóm xanh}{Rob Kolstad, 2000; dịch giả nguồn: BirDOR}

\iosec{Tóm tắt bài toán.}
John mua \(N\) con bò từ Blackbeard, và các con bò này thù ghét nhau. Ông định
nhốt chúng trong một khu đất kích thước \(X\times Y\). Vì thù ghét nhau, mỗi
con muốn càng xa các con khác càng tốt; chúng di chuyển cho tới khi tổng
khoảng cách Euclid giữa mọi cặp bò đạt cực đại.

Đây là bài nộp đáp án: hãy in một cách đặt vị trí. Bộ chấm sẽ so sánh tổng
khoảng cách của đáp án với đáp án chuẩn để tính điểm.

\iosec{Dữ liệu vào.}
Ba số nguyên \(N,X,Y\).

\iosec{Dữ liệu ra.}
In \(N\) dòng, mỗi dòng gồm hai tọa độ \(x_i,y_i\), với \(0\le x_i\le X\) và
\(0\le y_i\le Y\).

\iosec{Ví dụ.}
\begin{verbatim}
4 100 100
\end{verbatim}

\begin{verbatim}
100.00 0.00
0.00 100.00
100.00 100.00
0.00 0.00
\end{verbatim}

\problem{Cặp địa chỉ thân thiện}{amicbl}{Nhóm xanh}{Rob Kolstad, traditional; dịch giả nguồn: BirDOR}

\iosec{Tóm tắt bài toán.}
Sau khi chán nông trại, đàn bò chuyển ra ngoại ô. Mỗi địa chỉ là một số nguyên
dương. Hai địa chỉ được gọi là thân thiện nếu tổng các ước thực sự của địa chỉ
này bằng địa chỉ kia và ngược lại. Ví dụ, các ước thực sự của \(220\) có tổng
\(284\), còn các ước thực sự của \(284\) có tổng \(220\), nên \(220\) và
\(284\) là một cặp thân thiện.

\iosec{Dữ liệu vào.}
Hai số nguyên \(L,H\), với \(1<L<H<900000\).

\iosec{Dữ liệu ra.}
In mọi cặp địa chỉ thân thiện trong đoạn \((L,H]\), mỗi dòng một cặp, số nhỏ
in trước. Các dòng được sắp tăng theo số đầu tiên. Nếu không có cặp nào, tệp
ra có \(0\) dòng.

\iosec{Ví dụ.}
\begin{verbatim}
200 300
\end{verbatim}

\begin{verbatim}
220 284
\end{verbatim}

\section{USACO 2001 Spring}

\problem{Những con bò đốm}{spots}{Nhóm xanh}{Brian Dean, 2001}

\iosec{Tóm tắt bài toán.}
Farmer John muốn đưa \(N\) con bò ra hội chợ. Mỗi con bò đứng trong một ô
\(1\times 1\) trên một mặt phẳng lưới lớn, và con bò thứ \(i\) có tọa độ
\((r_i,c_i)\) cùng số đốm \(s_i\).

John được chọn một hình chữ nhật gồm \(A\) hàng và \(B\) cột, các cạnh song
song với trục hàng và cột, rồi đưa toàn bộ bò trong hình chữ nhật đó đi triển
lãm. Độ khác biệt của một hình chữ nhật là trị tuyệt đối của hiệu giữa số đốm
lớn nhất và nhỏ nhất của các con bò nằm trong đó. Hãy tìm độ khác biệt lớn
nhất trên mọi hình chữ nhật có thể chọn.

\iosec{Dữ liệu vào.}
Dòng đầu gồm \(N,A,B\). \(N\) dòng tiếp theo, mỗi dòng gồm \(r_i,c_i,s_i\), là
hàng, cột và số đốm của một con bò.

\iosec{Dữ liệu ra.}
In độ khác biệt lớn nhất.

\iosec{Ví dụ.}
\begin{verbatim}
8 4 2
1 4 9
1 5 8
2 10 2
3 2 6
4 6 1
5 15 3
6 4 5
7 9 4
\end{verbatim}

\begin{verbatim}
7
\end{verbatim}

\problem{Biểu thức bò}{cowfix}{Nhóm xanh}{Brian Dean, 2001}

\iosec{Tóm tắt bài toán.}
Cùng một biểu thức có thể viết ở dạng trung tố, tiền tố hoặc hậu tố. Tiền tố
và hậu tố không cần và không chấp nhận dấu ngoặc. Đàn bò của John viết
``biểu thức bò'', trộn lẫn các quy tắc trung tố, tiền tố và hậu tố. Biểu thức
chỉ gồm chữ số đơn và các toán tử nhị phân \(+\), \(-\).

Một biểu thức bò hợp lệ có thể được diễn giải theo nhiều cách. Hãy đếm số giá
trị khác nhau mà biểu thức đó có thể nhận.

\iosec{Dữ liệu vào.}
Một xâu không quá \(80\) ký tự, biểu diễn một biểu thức bò hợp lệ. Xâu không
có dấu cách.

\iosec{Dữ liệu ra.}
In số lượng giá trị khác nhau có thể thu được.

\iosec{Ví dụ.}
\begin{verbatim}
+54+-82
\end{verbatim}

\begin{verbatim}
2
\end{verbatim}

\problem{Tuyến vắt sữa tốt nhất}{route}{Nhóm xanh}{Brian Dean, 2001}

\iosec{Tóm tắt bài toán.}
Mỗi sáng, Farmer John cần lấy đủ \(K\) gallon sữa từ \(N\) con bò trên mặt
phẳng. Mỗi lần gặp một con bò, ông có thể lấy \(1\) gallon; nếu rời một con
bò, gặp con bò khác, rồi quay lại con bò cũ thì con bò cũ có thể lại cho thêm
\(1\) gallon.

Trên mặt phẳng có \(F\) hàng rào, mỗi hàng rào là một đoạn thẳng, không cắt
hàng rào khác và không đi qua bò. Hàng rào có chiều cao \(h\); xác suất John
vượt qua thành công hàng rào đó là \(1/h\). John luôn đi theo đoạn thẳng giữa
hai con bò, hoặc giữa chuồng và một con bò. Nếu đoạn đi trùng với hàng rào
hoặc đi qua đầu mút hàng rào, ông không phải vượt qua hàng rào đó. Xác suất
thành công của cả tuyến là tích xác suất khi vượt qua các hàng rào gặp trên
đường.

Tuyến vắt sữa bắt đầu ở chuồng, đi qua các bò để lấy đủ \(K\) gallon, rồi quay
về chuồng. Hãy tìm xác suất thành công lớn nhất.

\iosec{Dữ liệu vào.}
Dòng đầu gồm \(N,K,F,x,y\), trong đó \((x,y)\) là tọa độ chuồng. \(N\) dòng
tiếp theo là tọa độ các con bò. \(F\) dòng cuối, mỗi dòng gồm
\(x_1,y_1,x_2,y_2,h\), mô tả một hàng rào nối \((x_1,y_1)\) với
\((x_2,y_2)\) và có chiều cao \(h\).

\iosec{Dữ liệu ra.}
In xác suất thành công lớn nhất.

\iosec{Ví dụ.}
\begin{verbatim}
3 6 4 1 3
8 2
15 2
8 7
4 1 4 8 4
7 4 9 4 3
10 1 10 3 2
11 4 16 4 6
\end{verbatim}

\begin{verbatim}
1.3021e-03
\end{verbatim}

\problem{Những chuyến đi giữa hai chuồng}{barn}{Nhóm xanh}{Brian Dean, 1999}

\iosec{Tóm tắt bài toán.}
Có \(N\) mảnh đất đánh số từ \(1\) đến \(N\), nối với nhau bằng \(P\) lối đi
hai chiều. Mảnh đất \(1\) và \(2\) đều có chuồng, và giữa chúng không có lối đi
trực tiếp. Bessie muốn đi qua lại giữa hai chuồng càng nhiều lần càng tốt,
nhưng ngoài hai mảnh đất \(1\) và \(2\), cô không muốn đi qua cùng một mảnh đất
hai lần trong toàn bộ các chuyến đi. Hãy tính số chuyến đi tối đa có thể thực
hiện.

\iosec{Dữ liệu vào.}
Dòng đầu gồm \(N,P\). \(P\) dòng tiếp theo, mỗi dòng gồm hai mảnh đất được nối
bởi một lối đi.

\iosec{Dữ liệu ra.}
In số chuyến đi tối đa.

\iosec{Ví dụ.}
\begin{verbatim}
7 10
1 3
1 5
1 6
2 4
2 5
2 7
3 4
3 5
5 7
6 7
\end{verbatim}

\begin{verbatim}
3
\end{verbatim}

\section{USACO 2001 Open}

\problem{Tiếp sức bò}{relay}{Nhóm xanh}{Brian Dean, 2001}

\iosec{Tóm tắt bài toán.}
John chọn \(K\) con bò làm đội tiếp sức. Cuộc đua diễn ra trên \(N\) nông trại
đánh số từ \(1\) đến \(N\), nối bởi \(M\) con đường hai chiều; giữa hai nông
trại khác nhau có nhiều nhất một con đường. Mỗi đường có thời gian chạy qua
riêng.

Từng con bò lần lượt chạy từ nông trại \(1\) đến nông trại \(N\); con kế tiếp
chỉ bắt đầu sau khi con trước đã về đích. Theo luật, không hai con bò nào được
chạy đúng cùng một tuyến, trong đó tuyến là dãy nông trại đi qua. Một con bò
vẫn có thể đi qua cùng một nông trại nhiều lần. Hãy tính tổng thời gian nhỏ
nhất để cả đội hoàn thành.

\iosec{Dữ liệu vào.}
Dòng đầu gồm \(K,N,M\). \(M\) dòng tiếp theo, mỗi dòng gồm hai đầu mút của một
con đường và thời gian cần để một con bò đi qua đường đó.

\iosec{Dữ liệu ra.}
In tổng thời gian nhỏ nhất.

\iosec{Ví dụ.}
\begin{verbatim}
4 5 8
1 2 1
1 3 2
1 4 2
2 3 2
2 5 3
3 4 3
3 5 4
4 5 6
\end{verbatim}

\begin{verbatim}
23
\end{verbatim}

\section{Phạm vi còn lại}

Phần thân sau mục lục gồm phát biểu bài, dữ liệu vào, dữ liệu ra, ví dụ và
thông tin nguồn cho từng bài. Lớp văn bản PDF vẫn bị mojibake, nên các bài
chưa xuất hiện trong phần khôi phục ở trên cần tiếp tục được đọc từ ảnh render
và OCR. Các mục đã khôi phục được thêm theo từng đợt để giữ mẫu vào/ra và công
thức đủ chắc trước khi đưa vào bản LaTeX.

\end{document}
